\chapter{绪论}

\section{研究背景与意义}
\subsection{行业现状分析}
人工智能辅助编程正在从“代码补全工具”演进为“开发流程工具”。早期产品主要提供局部建议，开发者仍要手动完成上下文整理、代码拼接和运行验证。近两年出现的系统把需求输入、代码生成、结果校验和交付输出放在同一链路中，交互方式更接近任务驱动。行业关注点也因此从“能不能补全代码”转为“能不能稳定完成一个开发任务”。真实仓库评测结果说明，模型本体能力固然重要，但任务完成率同样受到流程编排、上下文注入和错误恢复机制影响（SWE-bench: Can Language Models Resolve Real-World GitHub Issues?）\cite{jimenez2024swebench}。

前端场景对这类变化更敏感。组件开发通常围绕模板结构、状态逻辑、事件交互和样式规则展开，需求描述天然具有较强结构性，适合自然语言输入。另一方面，前端交付强调“可见即所得”，如果系统只返回一段代码文本，用户无法快速确认视觉结果和交互行为，就很难把工具真正纳入日常工作流。这也是很多仅有“生成窗口”的产品在持续使用阶段遇到阻力的直接原因。

行业产品路线目前呈现出不同侧重点。在线生成平台把主要精力放在缩短需求到页面的距离，强调即时反馈和视觉交互（Introducing the new v0）\cite{vercel2026newv0}。代理式系统更关注多步骤任务执行，强调模型与工程环境的连接能力（OpenHands: AI-Driven Development）\cite{openhands2026github}。这两类路线共同推动了 AI 编程工具的工程化，但也让“稳定性、可追溯、安全边界”成为更明确的研究问题。

从教学项目和中小型团队的真实需求看，用户并不要求系统一次覆盖全部软件工程任务，而更关心单一任务链路是否足够稳定。自然语言到 UI 组件生成是一个边界清晰、可测性较高、结果可视化明显的切入点，适合作为 AI 辅助编程工具的工程化研究对象。

\begin{figure}[htbp]
\centering
\fbox{
\begin{minipage}{0.92\textwidth}
\centering
自然语言需求输入 $\rightarrow$ 任务创建与上下文绑定 $\rightarrow$ 流式生成与中间反馈 $\rightarrow$ 结构化抽取与安全校验 $\rightarrow$ 在线预览与版本下载
\end{minipage}}
\caption{AI 辅助前端开发的闭环路径}
\label{fig:chapter1-closed-loop}
\end{figure}

图 \ref{fig:chapter1-closed-loop} 展示了当前系统化工具常见的目标路径。闭环路径带来的直接收益是反馈更及时、错误更容易定位、结果更容易复用。

\subsection{现有系统存在的问题}
虽然行业工具数量增长较快，但把这些工具放进实际开发流程后，问题仍然集中在“链路完整性”上。用户在输入需求后，经常会遇到模型理解偏差。偏差不仅体现在界面布局，还体现在数据结构和交互规则的错配，导致结果看起来可用但难以直接进入项目。

输出稳定性不足也是常见问题。同一提示词在不同轮次生成的组件结构可能差异较大，脚本状态命名和样式细节变化频繁，给后续维护带来额外成本。对教学和课程项目而言，结果不可复现会直接影响实验记录和过程评分。

另一个明显问题是工程链路断点。很多工具把“生成”作为终点，没有持续提供在线预览、版本保存、历史回看和文件下载。开发者因此需要在多个平台之间反复复制内容，错误定位和成果留痕都不理想。对协作场景来说，缺少统一任务标识和版本编号也会影响沟通效率。

安全边界问题同样不能忽视。生成代码如果直接进入运行环境，可能引入外链脚本、动态执行或非预期网络访问。安全开发框架强调，风险控制应前移到开发和交付前阶段，而不是上线后再做补救（Secure Software Development Framework (SSDF) Version 1.1）\cite{nistssdf2024}。这类建议对 AI 代码生成系统同样适用。

\begin{table}[htbp]
\centering
\caption{自然语言到组件生成链路中的典型问题}
\label{tab:chapter1-problem-flow}
\small
\begin{tabular}{@{}p{0.18\textwidth}p{0.28\textwidth}p{0.34\textwidth}@{}}
\toprule
链路阶段 & 典型问题 & 对项目交付的影响 \\
\midrule
需求输入 & 描述歧义、约束不完整 & 生成结果偏离业务目标 \\
代码生成 & 输出结构波动、内容缺段 & 预览失败或需人工补全 \\
运行预览 & 缺少沙箱隔离与兜底 & 错误定位成本上升 \\
结果交付 & 无版本与下载管理 & 难以纳入日常开发流程 \\
\bottomrule
\end{tabular}
\end{table}

\begin{table}[htbp]
\centering
\caption{现有方案在工程闭环能力上的常见差异}
\label{tab:chapter1-gap}
\small
\begin{tabular}{@{}p{0.18\textwidth}p{0.14\textwidth}p{0.22\textwidth}p{0.14\textwidth}p{0.18\textwidth}@{}}
\toprule
方案类型 & 代码生成 & 在线预览 & 版本回溯 & 安全拦截 \\
\midrule
IDE 辅助型 & 较强 & 通常依赖本地工程 & 弱 & 依赖外部工具链 \\
在线生成型 & 较强 & 较强 & 中等 & 机制不统一 \\
代理执行型 & 中到强 & 依实现而定 & 中到强 & 依实现而定 \\
本文系统目标 & 强 & 强 & 强 & 前后端双层 \\
\bottomrule
\end{tabular}
\end{table}

表 \ref{tab:chapter1-gap} 表明本文课题的重点并不是重复实现“代码生成”本身，而是补足闭环中的薄弱环节，尤其是可预览、可追溯和可控执行三部分。

\subsection{本系统研发的必要性}
本课题来源于明确且持续的使用需求。用户希望通过自然语言描述直接得到可运行的 Vue 3 组件，并在同一页面完成预览和下载。这个目标看似聚焦在“生成代码”，实际上要求系统同时具备任务管理、状态反馈、版本存储和安全校验等工程能力。缺少其中任一环节，系统就会退化为一次性演示工具。

从实现成本角度看，课程设计与中小型团队普遍面临人力有限、迭代周期短的问题。组件搭建属于高频重复工作，自动化生成可以减少手工搭建时间，让开发者把精力投入在业务规则与交互细化。与此同时，教学场景强调过程可追踪，如果系统提供任务历史和版本记录，实验复盘和结果比对会更清晰。

从风险控制角度看，研发此类系统也有现实必要。模型输出不稳定是客观存在的，只有在工程层加入结构化抽取、异常处理和安全拦截，才能把不确定性限制在可控范围内。该思路与近年来软件工程研究强调的“系统能力与模型能力并重”方向一致（SWE-agent: Agent-Computer Interfaces Enable Automated Software Engineering）\cite{yang2024sweagent}。

\section{国内外研究现状}
\subsection{相关领域技术发展综述}
相关研究近几年变化很快。早期工作聚焦模型本身的生成能力，评价指标多为代码片段正确率。近期研究更重视真实工程任务完成率，评测对象也从静态题目逐步扩展到仓库级问题。这个变化让系统设计再次成为关键变量，提示词构造、上下文组织、执行接口和错误恢复会共同影响最终表现（SWE-bench: Can Language Models Resolve Real-World GitHub Issues?）\cite{jimenez2024swebench}。

代理式软件工程方向进一步说明，模型与环境之间的接口约束决定了任务是否可落地。即便模型生成能力较强，如果不能稳定读写工程上下文或处理运行反馈，也很难完成多步骤任务（SWE-agent: Agent-Computer Interfaces Enable Automated Software Engineering）\cite{yang2024sweagent}。这一结论对本文课题具有直接启发，即系统应把“可执行链路”作为核心对象，而不是把生成结果当作唯一目标。

在 UI 代码生成方向，近年研究把注意力放在结构一致性和视觉一致性。相关基准结果显示，模型在布局细节、组件层级和样式还原上仍存在稳定误差，这意味着后处理规则、格式约束和预览校验仍不可省略（Design2Code: How Far Are We From Automating Front-End Engineering?）\cite{si2024design2code}。因此，工程实现中通常需要“模型生成 + 规则修正 + 运行验证”的组合策略。

结合这些研究趋势可以看到，学术界与产业界的关注点正在收敛。产业侧重交付效率，学术侧重可解释评估，但两者都在强调系统化能力的重要性。本文系统的研究方向与这一收敛趋势保持一致，即在可控任务边界内提升完整链路质量。

\subsection{典型系统对比分析}
对比现有系统可以看到，不同产品的优化目标并不一致。在线页面生成平台通常强调输入到输出的速度，适合原型阶段和需求讨论阶段（Introducing the new v0）\cite{vercel2026newv0}。代理框架强调任务扩展能力，能够在更复杂的工程环境中执行多步操作（OpenHands: AI-Driven Development）\cite{openhands2026github}。两种路线都具有实际价值，但也分别面对可控性或复杂度方面的约束。

如果将这些路线直接迁移到本课题，会出现匹配度问题。单纯追求“快速出图”无法满足课程项目对可追溯和可下载的要求；直接引入通用代理框架又会增加部署和使用门槛。本文系统因此采用了边界更清晰的折中方案，把任务范围限定在“自然语言到 Vue 组件”，再在该范围内强化流式反馈、版本管理和安全控制。

这种定位使系统能够在有限开发周期内完成高价值闭环。用户获得的是一个可运行、可追踪、可交付的组件生成流程，而不是只看效果的演示页面。对于教学与中小规模开发场景，这种取舍更符合实际使用方式。

\begin{table}[htbp]
\centering
\caption{典型系统路线与本文系统定位对比}
\label{tab:chapter1-route}
\small
\begin{tabular}{@{}p{0.12\textwidth}p{0.22\textwidth}p{0.22\textwidth}p{0.24\textwidth}@{}}
\toprule
维度 & 在线生成平台路线 & 代理执行路线 & 本文系统路线 \\
\midrule
目标范围 & 快速产出页面 & 执行多步工程任务 & 稳定完成组件生成闭环 \\
交互重点 & 即时反馈与可视化 & 任务编排与工具调用 & 流式反馈与可追溯版本 \\
工程门槛 & 低到中 & 中到高 & 中 \\
适配场景 & 原型与演示 & 自动化工程任务 & 教学与中小型项目开发 \\
\bottomrule
\end{tabular}
\end{table}

表 \ref{tab:chapter1-route} 展示了本文系统的定位依据。该定位强调在可控范围内提高可用性，而不是一开始追求最大任务覆盖。

\section{研究目标与内容}
\subsection{系统核心功能定位}
本系统定位为“基于自然语言描述的 Vue 组件代码生成工具”，核心目标是让用户在浏览器内完成从需求输入到结果交付的完整流程。系统输出不是任意代码片段，而是包含 \texttt{template}、\texttt{script} 和 \texttt{style} 的可运行单文件组件。用户在生成过程中可以看到流式内容和任务状态，生成完成后可以继续完成预览、历史查询和文件下载。

围绕这一定位，系统功能组织遵循“闭环可验证”原则。输出层要求结构完整且可渲染，过程层要求状态可见且可追踪，安全层要求高风险模式在进入预览前被识别和拦截。该定位使系统具备工程使用价值，而不仅是展示模型能力。

\begin{figure}[htbp]
\centering
\fbox{
\begin{minipage}{0.92\textwidth}
\centering
研究目标 $\rightarrow$ 任务创建与流式生成 $\rightarrow$ 结构化抽取与安全校验 $\rightarrow$ 版本化存储与历史检索 $\rightarrow$ 在线预览与下载交付
\end{minipage}}
\caption{研究目标到系统能力的映射关系}
\label{fig:chapter1-goal-map}
\end{figure}

\subsection{拟解决的关键问题}
本课题实施中需要解决的问题集中在四个方面。模型输出格式并不总是稳定，系统必须保证在缺段或格式异常时仍可进入后续流程。流式生成涉及长连接与异步执行，任务状态转换需要保持一致，避免重复执行和状态覆盖。匿名访问模式下，系统还要保证不同用户数据隔离，同时保留版本回溯能力。生成代码进入预览前，则必须完成风险识别，避免不安全调用进入运行环境。

围绕这些问题，项目分别采用了结构化抽取与兜底补全、任务状态机与并发互斥、工作区键隔离、双层安全扫描等策略。每项策略都对应到后端服务和前端处理逻辑，便于单元测试和集成测试。通过这一组机制，系统将模型输出的不确定性压缩在可控范围内。

从研究角度看，上述问题并非实现细节，而是决定系统能否长期使用的基础条件。若只关注“能生成一次可见结果”，系统很难支撑持续迭代；若把稳定性、可追溯和安全边界作为同等目标，系统才具备工程落地价值。

\section{论文组织结构}
\subsection{章节内容概要说明}
全文共七章，内容安排与系统研发流程保持一致。第一章说明课题来源、研究价值与研究边界，明确本文关注的核心问题。第二章围绕系统落地所需技术给出选型依据，并结合代码实现说明关键机制。第三章基于实际业务场景展开需求分析，明确功能范围和非功能指标。第四章给出总体设计方案，涵盖系统架构、模块划分、数据库设计和安全控制策略。第五章聚焦关键实现过程，说明核心接口、算法处理和前后端协同逻辑。第六章展示测试方案和结果，对系统稳定性、可用性和安全性进行验证。第七章总结本文工作并给出后续可扩展方向。

该结构使论文叙述与工程开发路径一致，读者可以沿着“问题提出、技术选择、方案设计、实现验证”的逻辑理解全文内容，也便于将章节结论映射到实际代码实现。
